% ===================================================================
%  TABLO No 2 — İnsan vücudu. — Teneffüs. — Oyunlar.
%  Traduction 2026 d'apres texto/io, controlee sur texto/fr et
%  texto/hi. Consigne : voir texto/tr/10-tablo-01.tex.
% ===================================================================

%%K t02-apar-1 apar
\VUsaut{4.0mm}
\VUtitre{15.0pt}{60}{TABLO No 2}
\VUsaut{2.0mm}
\VUtitre{11.4pt}{0}{İnsan vücudu. --- Teneffüs.}
\VUtitre{11.4pt}{0}{Oyunlar.}

%%K t02-tit-0 sub
\VUtitre{13.2pt}{0}{İnsan vücudu.}
\VUsaut{2.4mm}
\VUcentre{10.2pt}{0}{\textit{(Basit bir doğa bilgisi dersi.)}}
\VUsaut{3.0mm}

%%K t02-01-1 p
1. --- İnsan vücudunun başlıca bölümleri şunlar: \VUgras{baş},
\VUgras{gövde} ve \VUgras{uzuvlar}.

%%K t02-tit-1 sub
\VUpk{10.2pt}{40}{Baş.}
\VUsaut{3.0mm}

%%K t02-02-1 p
2. --- Başın iki bölümü var: içinde \VUgras{beynin} bulunduğu ve
\VUgras{saçların} örttüğü \VUgras{kafatası}, bir de \VUgras{yüz}.

%%K t02-03-1 p
3. --- Resmimizde ne kafatası görünür ne beyin; ama yüzün başlıca bölümlerini
pek açık seçik görüyoruz.

%%K t02-04-1 p
4. --- Önce şu \VUgras{alın} var; güçlü bir zekâyı ve durmadan işleyen bir
kafayı belli eder. \VUgras{Şakaklar} çok açık. \VUgras{Göz} canlı ve ardına
kadar açık. Sık bir \VUgras{kaş} ile uzun \VUgras{kirpikler} onu korur.

%%K t02-05-1 p
5. --- Burada \VUgras{burun} ile \VUgras{burun delikleri} de var. Burunla
koklar ve soluk alırız. Burnun iki yanında \VUgras{yanaklar}, onların ötesinde
de \VUgras{kulaklar} bulunur. Yüzün alt bölümü \VUgras{ağız} ile
\VUgras{çeneden} oluşur.

%%K t02-06-1 p
6. --- Onları iyi göremiyoruz, çünkü \VUgras{bıyık} ile \VUgras{sakal} onları
gizler. Ama biliriz ki ağızda iki \VUgras{dudak}, \VUgras{dişleriyle} birlikte
\VUgras{üst çene} ile \VUgras{alt çene}, bir de \VUgras{dil} bulunur. Ağızla
konuşur ve yeriz. Dille yemeğimizin tadına bakarız.

%%K t02-07-1 p
7. --- Göz, kulak, burun ve ağız, parmaklarla birlikte, beş duyunun
uzuvları: \VUgras{görme}, \VUgras{işitme}, \VUgras{koklama},
\VUgras{tatma} ve \VUgras{dokunma}.

%%K t02-08-1 p
8. --- Böylece baş, insan vücudunun en ilgi çekici bölümlerinden biri.

%%K t02-tit-2 sub
\VUsaut{4.4mm}
\VUpk{10.2pt}{40}{Gövde.}
\VUsaut{3.0mm}

%%K t02-09-1 p
9. --- \VUgras{Boyun} başı gövdeye bağlar. \VUgras{Boğaz} ile \VUgras{ense}
ona girer.

%%K t02-10-1 p
10. --- Gövdede de pek çok gerekli organ bulunur. \VUgras{Göğüste}, soluk
aldığımız \VUgras{akciğerler} ile yaşamın merkezi olan \VUgras{yürek} var.
Yürek atmayı bıraktı mı, canlı ölür.

%%K t02-11-1 p
11. --- \VUgras{Kaburgalar} göğsü tutar.

%%K t02-12-1 p
12. --- Gövdenin öteki bölümleri \VUgras{böğür}, \VUgras{sırt},
\VUgras{omuzlar} ve \VUgras{karın} (göbek). Uyumak için yan yatar ya da
sırtüstü uzanırız, yükü de omuzda taşırız.

%%K t02-13-1 p
13. --- Karında \VUgras{mide} ile \VUgras{bağırsaklar} var. Yemeğimizi
sindirmemize yararlar.

%%K t02-tit-3 sub
\VUsaut{4.4mm}
\VUpk{10.2pt}{40}{Uzuvlar.}
\VUsaut{3.0mm}

%%K t02-14-1 p
14. --- Dört \VUgras{uzvumuz} var: iki \VUgras{kol} ile iki
\VUgras{bacak}. Kollarla nesneleri alır, tutar, kaldırır ve toplarız;
bacaklarla ayakta durur, yürür ve koşarız.

%%K t02-15-1 p
15. --- \VUgras{Omuz} kolu vücuda bağlar. Çok güçlü bir \VUgras{kas}, iki
başlı kas, kolu oyna; \VUgras{dirsek} de onu \VUgras{ön kola} bağlar.
\VUgras{Bilek}, \VUgras{parmaklarla} \VUgras{tırnaklarda} biten
\VUgras{elin} eklemini yapar. Elin üstünde, içinde \VUgras{kanın} aktığı bir
\VUgras{toplardamar} (bir de atardamar) görüyoruz.

%%K t02-16-1 p
16. --- Bacağın bölümleri şunlar: \VUgras{uyluk}, \VUgras{diz} ile
\VUgras{diz ardı}, \VUgras{etinin} \VUgras{deriyle} örtülü olduğu
\VUgras{baldır}, \VUgras{ayak bileği} ve \VUgras{ayak}. Bacağın
\VUgras{kemikleri} pek kalın ve pek sağlam. Ayağın ucunda \VUgras{ayak
parmakları} var. Öteki bölümleri \VUgras{taban} ile \VUgras{ökçe}.

%%K t02-17-1 p
17. --- İnsan vücudunun aşağı yukarı bütün açıklaması bu.

%%K t02-17-2 p
\textit{Not.} --- \textit{« ...\,m ağrıyor (başım, vb.) » cümlesinin yardımıyla
öğretmenler bütün bu sözleri iyi ya da kötü} \VUgras{beden sağlığı}
\textit{bakımından yeniden gözden geçirebilecekler.}

%%K t02-tit-4 sub
\VUsaut{5.0mm}
\VUpk{10.2pt}{40}{Jimnastik.}
\VUsaut{2.4mm}
\VUcentre{10.2pt}{0}{\textit{(Bir öğrencinin ağzından.)}}
\VUsaut{3.0mm}

%%K t02-18-1 p
18. --- Esnek, çevik ve sağlıklı olmak için bacaklarımızla kollarımızı
kullanmamız gerekir. Bunun için oyun oynar ve \VUgras{jimnastik} yaparız.

%%K t02-19-1 p
19. --- Hafta boyunca bu iki alıştırma okulun avlusunda olur (bakınız tablo
No 1). Ama perşembe ile pazar günleri koşacak ve eğlenecek yer bulduğumuz büyük
bir alana gideriz.

%%K t02-20-1 p
20. --- Öğretmenlerimiz ile anne babalarımız kimi zaman bize katılır; ama
alıştırmaları \VUgras{jimnastik öğretmeni} \textsuperscript{(12)} yaptırır.

%%K t02-21-1 p
21. --- Alanda, girişin solunda, \VUgras{jimnastik aletleri}
\textsuperscript{(1)} durur: \VUgras{merdivene} \textsuperscript{(2)}
tırmanırız, \VUgras{halkalarda} \textsuperscript{(3)} ve \VUgras{barfikste}
\textsuperscript{(4)} baş aşağı sallanırız, \VUgras{ip merdivenin}
\textsuperscript{(5)} ya da \VUgras{düğümlü ipin} \textsuperscript{(6)} ucuna
kollarımızla çıkarız.

%%K t02-22-1 p
22. --- Bu arada arkadaşlarımız \VUgras{tahta atın} \textsuperscript{(7)} ya
da \VUgras{paralel barların} \textsuperscript{(11)} üstünden atlar. Kimi
\VUgras{boks} \textsuperscript{(8)} yapar ya da maske ile \VUgras{flöret}
alıp \VUgras{eskrim} \textsuperscript{(9)} yapar. Kimi de sonunda
\VUgras{dambıl} \textsuperscript{(10)} kaldırır.

%%K t02-tit-5 sub
\VUsaut{4.6mm}
\VUpk{10.2pt}{40}{Oyunlar.}
\VUsaut{3.0mm}

%%K t02-23-1 p
23. --- Jimnastik yapanların çevresinde oğlanlarla kızlar türlü türlü
\VUgras{oyun} oynuyor. Kızlar \VUgras{çemberleriyle} \textsuperscript{(14)}
eğleniyor; \VUgras{ip atlıyorlar} \textsuperscript{(13)}, ya da erkek
kardeşleriyle kuzenlerini \VUgras{kroket} \textsuperscript{(15)} oyununa
çağırıyorlar. Karşı taraf çemberden kolayca geçeceğini sandığı anda, tahta
\VUgras{tokmaklarıyla} onun \VUgras{topunu} uzağa vuruyorlar, ve buna pek
seviniyorlar!

%%K t02-24-1 p
24. --- Oğlanlar daha güçlü oyunları sever. Sevine sevine \VUgras{dokuz kuka}
\textsuperscript{(16)} oynuyorlar; kukaları \VUgras{gülleyle}
\textsuperscript{(17)} temiz deviriyorlar. \VUgras{Topacı}
\textsuperscript{(18)} ile \VUgras{bilbokeyi} \textsuperscript{(19)}, hatta
\VUgras{kurbağa oyununu} \textsuperscript{(20)} de küçümsemezler. Ama en
sevdikleri oyunlar \VUgras{misket} \textsuperscript{(21)} ile \VUgras{duvar
topudur} \textsuperscript{(22)}, çünkü \VUgras{limonluğun}
\textsuperscript{(26)} duvarı hafif topu sektirmeye tam gelir.

%%K t02-25-1 p
25. --- Küçük Jan \VUgras{sırıklarının} \textsuperscript{(24)} üstünde pek
becerikli ve dengesini iyi tutuyor. Yanında birkaç arkadaşı aynı limonlukta ve
\VUgras{çitin} arkasında \VUgras{saklambaç} \textsuperscript{(25)} oynuyor.
Kimi \VUgras{el oyununu} \textsuperscript{(28)} daha çok seviyor, ya da bir
\VUgras{sopadan} ötekine uçup duran \VUgras{kuşu} \textsuperscript{(29)}.

%%K t02-noto-1 noto
\VUnotes{143.2mm}{%
(*) Çocuk oyunlarının adları çoğu kez büsbütün ulusal. Onlara idoda
elimizden geldiğince ad verdik — kimi zaman kendilerini belli eden eylemden,
kimi zaman da pek yanıltıcı olmadıkça şu ya da bu ülkenin adını alarak.
Söz gelişi: \VUgras{fıçı oyunu} (almanca ile fransızca) ile
\VUgras{kurbağa oyunu} \textsuperscript{(20)} (ingilizce) aynı; oyuncular
yuvarlak pullarını, delikleri olan sandığın (fıçının) üstünde ağzı açık duran
madenî kurbağanın ağzına atmaya çalışır. \VUgras{Omuz atlaması}
\textsuperscript{(30)} ile \VUgras{keçi atlaması} arasındaki ayrım yalnız
şu: başın üstünden atlamak için oyuncular ellerini arkadaşlarının
omuzlarına koyar, oysa « \VUgras{keçi atlamasında} » yalnız sırtına koyarlar.
--- Ya da kimi oyuncular eğilir ve ötekiler ellerini onların sırtına koyup
üstlerinden atlar; öncekilerin kımıldamadan itilmeye dayanması, sonrakilerin
de dengelerini yitirmeden atlaması gerekir (tablonun kendisine bakınız).%
}

%%K t02-26-1 p
26. --- Alanın ortasında iki ağaç var. Bunlardan birinin altında yedi sekiz
oğlan \VUgras{omuz atlaması} \textsuperscript{(30)} (*) oyununu kurmuş. Bu
oyun her ülkede bilinmez; ama \VUgras{keçi atlamasının} ne olduğunu herkes
bilir, ki çocuklar şu sırada onu oynamıyor.

%%K t02-tit-6 sub
\VUsaut{5.0mm}
\VUpk{10.2pt}{40}{Açık hava oyunları.}
\VUsaut{3.4mm}

%%K t02-27-1 p
27. --- \VUgras{Körebe} \textsuperscript{(31)} de pek eğlenceli; kızlarla
oğlanlar sırayla gözlerini \VUgras{bağlar} ve yakalanmaya bırakan
beceriksizleri yakalarlar.

%%K t02-28-1 p
28. --- Alanın ortasında, işte biraz kaba olsa da pek fransız bir oyun:
\VUgras{ayı} \textsuperscript{(32)}, ki \VUgras{uçurtmanın}
\textsuperscript{(33)} sakin oyunundan, hatta ingiliz oyunu \VUgras{tenisten}
\textsuperscript{(34)} de, çok daha gürültülü.

%%K t02-29-1 p
29. --- Bu son oyun büyük öğrencilerin sevinci. Öğretmenlerimiz bile onu
seve seve benimser. Çok beceri ile çeviklik ister, ve bana pek hoş gelir.
\VUgras{Salıncağı} \textsuperscript{(35)} kızlara bırakırım.

%%K t02-30-1 p
30. --- Kolayca tehlikeli olabilen bir başka ingiliz oyunu da hoşuma gitmez.

%%K t02-30-1b p
\VUgras{Futbolu} \textsuperscript{(36)} demek istiyorum, ki ağabeyimin
sevinci. Bana eski oyunumuz \VUgras{esir almaca} \textsuperscript{(38)} daha
hoş gelir, o kadar sağlıklı ve çok daha az sert.

%%K t02-tit-7 sub
\VUsaut{4.6mm}
\VUpk{10.2pt}{40}{Bisiklet ve top oyunları.}
\VUsaut{3.0mm}

%%K t02-31-1 p
31. --- Alanın çevresinde \VUgras{bisiklete} \textsuperscript{(37)} binmek de
hoşuma gider.

%%K t02-32-1 p
32. --- Gelecek yıl \VUgras{binicilik} \textsuperscript{(39)} öğreneceğim. At
sekerek ya da düzgün tırısla giderken, dörtnala engelleri aşarken ne güzel
görünür!

%%K t02-33-1 p
33. --- Yazın \VUgras{yüzme} \textsuperscript{(40)} öğreniriz. Irmağın duru
soğuk suyunda seve seve dalar ve yıkanırız. Kimi zaman sonunda bir kâğıt
\VUgras{balon} \textsuperscript{(41)} uçururuz, ki sıcak havayla dolar dolmaz
ağır ağır yükselir.

%%K t02-34-1 p
34. --- Daha pek çok oyun var, ama onları saya saya hiç bitiremem, ve bu
küçük anlatıdan da görünür ki güzel alanımızda perşembeler tatsız olmaktan çok
uzak.

%%K t02-34-2 p
N. B. --- \textit{Öğretmenler her önemli oyunu daha ayrıntılı anlatarak bu
bölümü kolayca tamamlayabilecekler.}
\VUsaut{6.0mm}
\VUfilet{22mm}
